\documentclass[12pt, a4paper]{article}

%%------------------Preamble------------------%%

    %% 引入Package %%
\usepackage[margin=2cm]{geometry} % 上下左右距離邊緣3cm
\usepackage{fancyhdr}   % 頁首頁尾 fancy header
\usepackage{titling}    % 設定 title,author,date
\usepackage{mathtools,amsthm,amssymb} % 引入 AMS 數學環境
\usepackage[shortlabels,inline]{enumitem} %加強 enumerate, itemize 功能
\usepackage{setspace}
\usepackage{tcolorbox}
\usepackage{xcolor} % 用於表格上色
\usepackage{graphicx}      % 用於縮放表格寬度
\usepackage{array}
\usepackage{float} % for [H] float placement
\tcbuselibrary{breakable}
\usepackage{tabularx}
\usepackage{colortbl}

    %% 中文環境，僅限使用 XeLaTeX 編譯器 %%
\usepackage[CJKmath=true,AutoFakeBold=3,AutoFakeSlant=.2]{xeCJK}

\setCJKmainfont{Noto Serif CJK TC}
\setCJKsansfont{Noto Sans CJK TC}
\setCJKmonofont{Noto Sans Mono CJK TC}

\newCJKfontfamily\Kai{Noto Serif CJK TC}
\newCJKfontfamily\Hei{Noto Sans CJK TC}
\newCJKfontfamily\NewMing{Noto Serif CJK TC}

\XeTeXlinebreaklocale "zh"




\begin{document}

\title{T7 選擇題解答} % 設定標題
\author{李原廷} % 設定作者
\date{\today} % 設定完成日期，也可以使用 \date{\today}

\maketitle



\begin{enumerate}

    \item 關於我國全民健康保險所採用的財務處理方式，下列敘述何者錯誤？
    \begin{enumerate}[label=\Alph*.]
        \item 會面臨長期性基金的投資運用問題
        \item 安全準備金的投資首重安全性與流動性
        \item 制度財務的穩定性極易受到人口老化的影響
        \item 較不易落實受益原則
    \end{enumerate}
    解答：
    \begin{enumerate}[label=\Alph*.]
        \item 採行的是隨收隨付制，當年度收到的保費，幾乎都在當年度支付給醫療院所了。留在手上的「安全準備金」規模非常小，且隨時可能因為傳染病爆發需要緊急提領付錢。因此，健保基金沒有餘裕進行「長期性基金的投資運用」（長期投資會把錢卡死，面臨流動性風險）。
        \item 健保的安全準備金主要目的是為了「救急與周轉」。因此，這筆錢在管理上以「安全性」與「流動性（隨時能變現）」第一位，收益性反而是最次要的考量。
        \item 隨收隨付制最大的致命傷就是極易受到人口結構變化的衝擊。人口老化意味著「繳納保費的年輕工作者變少」，同時「需要大量使用醫療資源的老年人變多」，這會直接導致健保收支嚴重失衡。
        \item 健保具有強烈的「所得重分配」與「社會互助」性質。保費是按「負擔能力（薪資高低）」來收，而不是按「生病機率或使用醫療資源的多寡」來收。因此，繳最多錢的人通常不是看最多次病的人，這導致健保制度不易落實「受益原則（誰受益、誰付費）」。
    \end{enumerate}

    \item 針對國內現行全民健康保險的「部分負擔」設計（除特殊情況外，保險對象就醫時，通常尚需要負擔法定金額或比例的醫藥費用），下列何者並非可能達到的效益？
    \begin{enumerate}[label=\Alph*.]
        \item 減輕制度的醫療支出負擔
        \item 降低保險對象的道德危機情況
        \item 降低制度的逆選擇情況
        \item 降低保險對象對醫療資源的不當利用情況
    \end{enumerate}
    解答：略。

    \newpage

    \item 有關就業保險之敘述，下列何者錯誤？
    \begin{enumerate}[label=\Alph*.]
        \item 就業保險法於民國92年1月1日正式施行
        \item 給付水準最高可領到平均月投保薪資的80\%
        \item 無論是自願或非自願離職均可請領
        \item 申請人須具有工作能力及繼續工作意願
    \end{enumerate}
    解答：
    就業保險中最主要的「失業給付」，其設立宗旨在於保障勞工因「非自願離職」（如：公司倒閉、裁員、資遣、因雇主違反勞基法而終止契約等）所造成的短期收入中斷。如果勞工是「自願離職」（如：個人生涯規劃、跳槽、單純不想做），無法請領失業給付。

    \item 有關公共年金財務規劃方式的確定給付制及確定提撥制，下列敘述何者正確？
    \begin{enumerate}[label=\Alph*.]
        \item 確定給付制屬於量入為出，確定提撥制則是量出為入的概念
        \item 確定給付制下，無須設置基金
        \item 我國勞工保險財務規劃，都屬於確定提撥制
        \item 確定給付制較容易因社會、經濟、人口結構變化等因素，造成年金制度財務不穩定
    \end{enumerate}
    解答：
    \begin{itemize}
        \item 確定給付制 (DB)：承諾未來退休時，依照你的「年資與薪資」發給一個「確定金額」的退休金（例如現在的勞保老年年金、軍公教退撫舊制）。
        \item 確定提撥制 (DC)：規定現在每個月要繳「確定比例」的保費存入你的個人專戶，未來退休時，專戶裡本金加投資收益有多少，就領多少（例如現在的勞工退休金新制）。
    \end{itemize}
    \begin{enumerate}[label=\Alph*.]
        \item 確定給付制 (DB) 是先決定未來要給你多少錢（出），再去推算現在每個人該收多少保費（入），屬於「量出為入」。確定提撥制 (DC) 是現在專戶存了多少錢（入），未來就只能領多少錢（出），屬於「量入為出」。
        \item 即使是確定給付制，通常也會採用「部分儲備制」或「完全儲備制」，必須設置龐大的年金基金（如我國的勞保基金、退撫基金）並進行投資運用，以因應未來的給付承諾。
        \item 「勞工保險（勞保）」是採用「確定給付制 (DB)」；而「勞工退休金（勞退新制個人專戶）」才是採用「確定提撥制 (DC)」。兩者制度不同。
        \item 在確定給付制 (DB) 下，政府或雇主承諾了未來的給付金額。一旦發生人口老化、少子化（繳保費的年輕人變少，領錢的老人變多且活更久）、或經濟不景氣（基金投資收益不佳），這個承諾的給付總額依然無法改變，最終就會導致整個年金制度的財務產生巨大缺口（如台灣目前的勞保破產危機）。
    \end{enumerate}


    \item 關於我國勞工退休金的提繳，下列敘述何者錯誤？
    \begin{enumerate}[label=\Alph*.]
        \item 勞工得自願提繳退休金
        \item 勞工自願提繳在其每月工資6\%範圍內，得自當年度個人綜合所得總額中全數扣除
        \item 雇主應為適用勞動基準法之勞工，按月提繳退休金儲存勞工退休金個人專戶
        \item 雇主、自營作業者、攤販等不可以參與提繳及請領退休金
    \end{enumerate}
    解答：

    勞工退休金條例
    \begin{enumerate}[label=\Alph*.]
        \item 第 14 條
        \item 第 14 條第 3 項
        \item 第 6 條與第 14 條
        \item 第 7 條第 2 項
    \end{enumerate}

    \item 就社會保險的財務處理方式而言，關於部分提存準備制之敘述，何者錯誤？
    \begin{enumerate}[label=\Alph*.]
        \item 會帶來一定程度的世代間所得重分配效果
        \item 基金或安全準備金金額，會高於隨收隨付制
        \item 基金或安全準備金在投資運用上，須考量長期收益
        \item 對國家總體經濟，不會帶來資本存量的影響
    \end{enumerate}
    解答：
    部分儲備制 (Partially Funded System)：
        \begin{itemize}
            \item 運作原理：如我國勞保的設計，當期保費收入大於支出時會累積準備金，但準備金規模又不足以完全涵蓋未來的潛在負債。
            \item 結論：雖然準備金不如完全儲備制龐大，但仍累積了相當規模的基金，因此仍然需要且重視長期投資收益來挹注財源、延緩破產。
        \end{itemize}

    \item 下列何者不屬於降低逆選擇問題的處理方式？
    \begin{enumerate}[label=\Alph*.]
        \item 政府訂定住宅租賃定型化契約應記載及不得記載事項，來保障租屋者的權益
        \item 政府針對勞工紓困貸款訂定相關紓困條件
        \item 政府制定食品安全衛生管理法，來維護國民健康
        \item 政府嚴格審核就業保險的失業給付請領條件
    \end{enumerate}
    解答：
    \begin{itemize}
        \item 逆選擇： 發生在「交易前（事前）」的資訊不對稱。因為買方不知道賣方產品的真實品質（或保險公司不知道投保人的真實風險），導致市場上只剩下「劣質品（檸檬車）」或「高風險者」。解決方法通常是強制揭露資訊、強制納保、提供品質認證等。
        \item 道德風險： 發生在「交易後（事後）」的資訊不對稱。因為簽約或獲得保障後，當事人的行為改變（變得不小心、不努力），且另一方無法輕易監督。解決方法通常是設立條件限制、部分負擔（自付額）或加強事後審核監督。
    \end{itemize}
     (D)：當勞工處於失業狀態並開始領取失業給付（交易後），他們可能會產生「既然有錢領，乾脆不積極找工作」的行為改變，政府又很難隨時監視他們有沒有在投履歷。為了解決這種「事後的道德風險」，政府才會嚴格審核請領條件（例如：要求每個月必須提供多次求職被拒絕的證明），這不是為了解決事前的逆選擇問題。

    \item 國內全民健康保險採行全民納保的方式，主要有助於減輕下列那個（些）問題？
    \begin{enumerate}[label=\Alph*.]
        \item 逆選擇
        \item 道德危機
        \item 逆選擇與道德危機
        \item 都無助於減輕逆選擇與道德危險問題
    \end{enumerate}
    解答：略。

    \item 有關於社會保險費的計算，下列敘述何者正確？
    \begin{enumerate}[label=\Alph*.]
        \item 社會保險的社會性特質，衍生出社會保險「繳費與給付不完全對等」的作法
        \item 社會保險按照性別、年齡、風險值來收費
        \item 全民健康保險會因為保費高低而有不同給付水準
        \item 社會保險保費通常由政府全額負擔
    \end{enumerate}
    解答：
    \begin{enumerate}[label=\Alph*.]
        \item 社會保險強調「社會連帶」與「所得重分配」的社會性特質。為了確保弱勢群體也能獲得基本保障，制度設計上會讓高所得者多繳費、低所得者少繳費，但大家生病時獲得的醫療保障卻是一樣的。這種財富重分配的機制，必然導致個人「繳交的保費」與「獲得的給付」不完全對等（不完全成比例）。
        \item 按照性別、年齡、身體健康狀況（風險值）來精算收費，是「商業保險」的作法。社會保險（如我國健保、勞保）通常是依據「負擔能力（如：薪資所得高低）」來收取保費，而非依據個人風險高低收費。
        \item 全民健康保險的宗旨是保障全體國民平等的就醫權利。只要納保，無論你每個月繳的是最高級距的保費還是最低級距的保費，生病時去醫院看診所享受的「醫療給付水準」是完全相同的，不會因為保費高低而有差別待遇。
        \item 社會保險的財源通常採用「共同負擔（三方負擔）」的原則，由被保險人（勞工）、雇主（投保單位）與政府依比例共同分擔保費，而不是由政府全額負擔（若全由政府負擔，那是社會福利津貼或社會救助，而非社會保險）。
    \end{enumerate}

    \item 社會保險財務處理方式中，基金（或安全準備金）的投資運用，下列何者較不涉及長期收益？ 1.隨收隨付制 2.完全儲備制 3.部分儲備制？
    \begin{enumerate}[label=\Alph*.]
        \item 2.
        \item 1.
        \item 3.
        \item 1.及2.
    \end{enumerate}
    解答：
    \begin{enumerate}[1.]
        \item 隨收隨付制 (Pay-As-You-Go, PAYGO)：
        \begin{itemize}
            \item 運作原理：以「當期」的保費收入，直接支付「當期」的保險給付（例如：現在年輕人繳的保費，直接拿去付給現在的老人當退休金）。
            \item 基金特徵：這種制度原則上不累積龐大的基金，只會保留非常少量的「安全準備金」，目的僅在於應付短期現金流的波動與周轉。
            \item 結論：因為準備金規模小且必須隨時提領，只能放在定存等流動性極高的短期資產，因此幾乎不涉及（也不依賴）長期投資收益。
        \end{itemize}
        \item 完全儲備制 (Fully Funded System)：
        \begin{itemize}
            \item 運作原理：每一代人（或每個人）現在繳納的保費，都會被完整存起來並進行投資，等未來自己發生保險事故（或退休）時，再拿出來支付給自己。
            \item 結論：這種制度會累積龐大的基金池，其財務穩健度依賴長期投資的複利收益，以對抗通貨膨脹並支應未來的給付。
        \end{itemize}
        \item 部分儲備制 (Partially Funded System)：
        \begin{itemize}
            \item 運作原理：介於上述兩者之間（如我國勞保的設計），當期保費收入大於支出時會累積準備金，但準備金規模又不足以完全涵蓋未來的潛在負債。
            \item 結論：雖然準備金不如完全儲備制龐大，但仍累積了相當規模的基金，因此仍然需要且重視長期投資收益來挹注財源、延緩破產。
        \end{itemize}
    \end{enumerate}


    \item 關於在社會保險及商業保險等二者中的逆選擇問題，下列敘述何者錯誤？
    \begin{enumerate}[label=\Alph*.]
        \item 社會保險的逆選擇發生率會高於商業保險
        \item 社會保險與商業保險的逆選擇情況都是起因於資訊不對稱
        \item 透過擴大強制納保對象可以降低社會保險的逆選擇問題
        \item 透過強化核保可以降低商業保險的逆選擇問題
    \end{enumerate}
    解答：
    社會保險（如我國的全民健保、勞保）為了解決逆選擇問題採用「強制納保（Compulsory Enrollment）」。強制納保把健康的人（低風險）和常生病的人（高風險）全部綁在同一個風險池裡，讓低風險者無法退出。因此，社會保險的逆選擇發生率實際上會「低於」自由投保的商業保險。

    \item 有關勤勞所得租稅抵減（earned income tax credit，EITC）與稅式支出（tax expenditure，TE）之敘述，下列何者正確？
    \begin{enumerate}[label=\Alph*.]
        \item EITC是給予低所得者所得減除的優惠
        \item EITC對低勞動所得者所形成之所得效果必使勞動供給增加
        \item TE可能形成顛倒式的補貼（upside-down subsidy）
        \item 若TE使民眾的支出金額大於政府直接支出，則其價格彈性絕對值小於1
    \end{enumerate}
    解答：
    \begin{enumerate}[label=\Alph*.]
        \item EITC是直接給予低所得勞工的「應納稅額扣抵（Tax Credit）」，甚至若扣抵後稅額為負，政府還會退還現金。它不是單純在計算所得時的「所得減除（Income Deduction）」。
        \item EITC 雖然會透過「替代效果」（提高實質工資）鼓勵低所得者進入職場或多工作，但就「所得效果」而言，因為 EITC 實質上增加了勞工的所得（變有錢了），在休閒為正常財的假設下，變有錢反而會促使人想多休閒、減少勞動供給。
        \item 稅式支出（Tax Expenditure）是指政府透過稅法上的特別規定（如：免稅額、扣除額、抵減額等），讓特定納稅人減少稅負，實質上等同於政府發放補貼。如果稅式支出是採用「列舉扣除額（Deductions）」的形式（例如：購屋借款利息扣除、捐贈扣除），其實質補貼利益會等於「扣除額 × 該納稅人的邊際稅率」。

        我國採累進稅率，高所得者的邊際稅率較高，低所得者的邊際稅率較低。這會導致同樣捐款 1 萬元，富人省下的稅（獲得的補貼）比窮人還要多，這種「越有錢拿越多補助」的現象就稱為顛倒式補貼（Upside-down Subsidy）。
        \item 稅式支出若視為一種價格補貼（降低購買特定物品的價格），如果要讓民眾增加的支出金額「大於」政府損失的稅收（即政府直接支出的等價物），這代表民眾對該物品的需求數量必須對價格變動非常敏感。也就是說，其需求的價格彈性絕對值必須大於 1（富於彈性），而非小於 1。
    \end{enumerate}

    \item 全民健康保險所提供的預防保健服務，下列敘述何者正確？
    \begin{enumerate}[label=\Alph*.]
        \item 可改善逆選擇問題
        \item 為避免浪費醫療資源，收取高額部分負擔
        \item 為現金補助措施之一
        \item 免除部分負擔
    \end{enumerate}
    解答：
    \begin{enumerate}[label=\Alph*.]
        \item 「逆選擇（Adverse Selection）」是指在保險市場中，高風險者（常生病的人）傾向投保，而低風險者（健康的人）不願投保的現象。我國健保是採用「強制納保」來解決逆選擇問題，而非透過預防保健服務來解決。
        \item 收取高額部分負擔會「阻礙」民眾進行健康檢查的意願，這與預防醫學的初衷不符。針對一般門診或急診才會利用部分負擔來抑制不必要的醫療浪費。
        \item 全民健保提供的預防保健是直接給予民眾免費的健康檢查與篩檢服務，這屬於「實物補助（In-kind Subsidy）」，而非直接發現金給民眾的「現金補助」。
        \item 預防保健具有「正向外部性」，能達到「早期發現、早期治療」的效果，從長遠來看反而能大幅降低健保治療重症的鉅額支出。為了鼓勵民眾積極參與預防保健，政府在政策設計上特地免除收取部分負擔，降低民眾的使用門檻。
    \end{enumerate}

    \item 社會救助法的中央主管機關為：
    \begin{enumerate}[label=\Alph*.]
        \item 內政部
        \item 衛生福利部
        \item 勞動部
        \item 財政部
    \end{enumerate}
    解答：

    第 3 條
    \begin{enumerate}[1.]
        \item 本法所稱主管機關：在中央為衛生福利部；在直轄市為直轄市政府；在縣（市）為縣（市）政府。
        \item 本法所定事項，涉及各目的事業主管機關職掌者，由各目的事業主管機關辦理。
    \end{enumerate}


    \item 由於我國人口結構發生變化，導致勞保財務問題日益嚴重，下列何者非主要原因？
    \begin{enumerate}[label=\Alph*.]
        \item 依賴人口比率日增
        \item 男女性別比例失衡
        \item 65 歲以上人口比例日漸增加
        \item 生育率逐年下滑
    \end{enumerate}
    解答：略。

    \item 有關李嘉圖等值定理（Ricardian equivalence theory），下列敘述何者正確？
    \begin{enumerate}[label=\Alph*.]
        \item 在既定的政府支出下，以發行公債或課稅方式融通，其經濟效果一致
        \item 在既定的政府支出下，以發行公債或印鈔票方式融通，其經濟效果一致
        \item 最適租稅應使所有財貨在課稅後，受補償需求量下降的幅度一樣
        \item 最適租稅應使所有財貨在課稅後，受補償需求量上升的幅度一樣
    \end{enumerate}
    解答：略。

    \item 下列何種預算制度最能顯示需要資金之目標，為達成這些目標之計畫成本，以及每一計畫所完成工作之數量資料？
    \begin{enumerate}[label=\Alph*.]
        \item 設計計畫預算
        \item 績效預算
        \item 零基預算
        \item 功能預算
    \end{enumerate}
    解答：
    \begin{enumerate}[label=\Alph*.]
        \item （PPBS）著重於「長程規劃」與「成本效益分析」，強調不同替代方案的選擇，而非單純的日常工作數量與單位成本。
        \item （ZBB）核心精神在於將預算建立在「政府究竟做了什麼事」之上，而不是單純列出「政府買了什麼東西」：要資金之目標（或功能）、達成目標之計畫成本、完成工作之數量資料（績效衡量）
        \item 著重於每年預算編製時「不考慮過去的既有基礎（歸零）」，所有計畫都要重新審視、由下而上提出決策包並進行優先順序排列。
        \item 僅是將預算按政府的大政方針（如國防、教育、社會福利）進行分類，缺乏詳細的計畫成本與具體工作數量的衡量。
    \end{enumerate}

    \item 關於新古典經濟學派對於財政赤字之看法，下列敘述何者正確？
    \begin{enumerate}[label=\Alph*.]
        \item 政府可以赤字預算調節景氣循環
        \item 發行公債將造成排擠效果
        \item 發行公債可使後代子孫擁有較多之資本資產
        \item 發行公債可提升經濟效率
    \end{enumerate}
    解答：
    \begin{enumerate}[label=\Alph*.]
        \item 新古典學派認為市場具備自我調節能力，反對政府過度干預；主張政府應以「赤字預算調節景氣循環（反景氣循環政策）」的是凱因斯學派。
        \item 新古典學派認為，在資源固定的情況下，政府為了彌補財政赤字而「發行公債」，等於是政府進入可貸資金市場與民間企業搶錢。這會導致市場對資金的需求增加，進而推升市場利率。利率上升後，民間企業的借貸成本變高，便會減少民間的投資意願。這種政府支出增加（發債）導致民間投資減少的現象，稱為「排擠效果（Crowding-out Effect）」。
        \item 發行公債會產生排擠效果，導致民間「投資減少」。既然現在的投資減少了，未來所能累積的資本存量自然也會跟著減少。因此，發行公債反而會使後代子孫擁有較少的資本資產。
        \item 新古典學派認為，舉債會排擠民間投資，且未來為了還債必須提高稅收，這會扭曲價格機能並產生超額負擔。因此，發行公債非但不能提升，反而會損害經濟效率。
    \end{enumerate}

    \item 目前有關公共債務法及預算法之規定，下列敘述何者正確？
    \begin{enumerate}[label=\Alph*.]
        \item 政府不得制定法律排除公共債務法之規定，但是得排除預算法之規定
        \item 政府不得制定法律排除預算法之規定，但是得排除公共債務法之規定
        \item 政府不得制定法律排除公共債務法及預算法之規定
        \item 政府得制定法律排除公共債務法及預算法之規定
    \end{enumerate}
    解答：
    財政紀律法第1條第2項明定：「有關財政紀律之規範，本法未規定者，適用其他有關法律之規定。」確立出該法優先於其他原有法律之特性。所以兩法皆得制定法律排除公共債務法及預算法之規定。

    \item 政府本年度總預算編列償還債務利息支出 2 億元，償還債務本金支出 14 億元，購買房屋支出 10 億元，增加投資支出 4 億元。根據我國預算法相關規定，資本支出金額總計為：
    \begin{enumerate}[label=\Alph*.]
        \item 10 億元
        \item 14 億元
        \item 28 億元
        \item 30 億元
    \end{enumerate}
    解答：
    \begin{itemize}
        \item 經常支出與資本支出的界定（歲出部分）：
        \begin{itemize}
            \item 經常支出：指維持政府日常運作所需要的經費，例如人事費、業務費、補助費，以及「債務利息支出」（ 2 億元屬於經常支出）。
            \item 資本支出：指用於增置、擴充、改良資產或增加政府投資的支出。「購買房屋支出 10 億元」（實體資產的增加）與「增加投資支出 4 億元」（金融資產的增加）均屬於資本支出。
        \end{itemize}
        \item 債務償還的特殊規定（非歲出）：
        \begin{itemize}
            \item 稱歲出者，謂一個會計年度之一切支出。但不包括債務之償還。
            \item 「償還債務本金支出」（14 億元）在預算編製上屬於「融資調度項目」（負債的減少），它既不屬於經常支出，也不屬於資本支出。
        \end{itemize}
    \end{itemize}
    資本支出總計 ＝ 購買房屋支出（10 億元） ＋ 增加投資支出（4 億元） ＝ 14 億元。

    \item 布坎南（Buchanan）認為公債是否存在代際間負擔移轉，須視下列何者而定：
    \begin{enumerate}[label=\Alph*.]
        \item 公債之收入用途
        \item 公債之發行方式
        \item 公債之發行主體
        \item 公債之發行對象
    \end{enumerate}
    解答：布坎南（James Buchanan）認為，公債是否造成代際間負擔移轉，關鍵在於公債收入用在什麼地方。
    \begin{itemize}
        \item 若公債收入用於當代消費性支出，例如一般行政支出、社會福利或非生產性支出，則當代人享受利益，未來世代負擔還本付息，會產生代際負擔移轉。
        \item 若公債收入用於公共投資，例如交通建設、教育、基礎設施等，未來世代雖要繳稅償債，但也享受公共投資帶來的利益，負擔可能被利益抵銷，較不構成不公平的代際負擔移轉。
    \end{itemize}
    因此，布坎南認為重點不只是政府有沒有發行公債，而是公債籌得的資金是否用於能使未來世代受益的用途。

    \item 何種預算制度可由公民投票決定計畫之優先順序？
    \begin{enumerate}[label=\Alph*.]
        \item 功能預算
        \item 績效預算
        \item 零基預算
        \item 參與式預算
    \end{enumerate}
    解答：
    \begin{enumerate}[label=\Alph*.]
        \item 依據政府的「政事功能」（如國防、教育、經濟發展）來進行預算分類編列，主要目的是讓外界清楚政府的錢花在哪一類的施政項目上，不涉及公民投票。
        \item 強調「投入與產出」的關係，將預算與施政計畫的「具體績效指標（成本與工作量）」結合，目的是提升行政效率與資源運用效能。
        \item 不考慮過去的預算額度（不沿用「歷史基期」加加減減的做法），每年編列預算時都必須「歸零」重新評估。所有計畫都必須重新審查其必要性與成本效益，以決定資源分配。
        \item 一種直接民主的展現。政府將一部分的公共預算資源交由在地居民，讓民眾透過公開討論、提案，最後由「公民投票」的方式來決定要執行哪些公共建設計畫以及其優先順序。
    \end{enumerate}

    \item 為強化債務管理，中央政府當年度稅課收入 1.5 兆元，國營事業繳庫收入 2,000 億元，至少應編列多少債務還本：
    \begin{enumerate}[label=\Alph*.]
        \item 150 億元
        \item 170 億元
        \item 750 億元
        \item 850 億元
    \end{enumerate}
    解答：

    根據《公共債務法》第 12 條第 1 項規定：「中央、直轄市、縣（市）及鄉（鎮、市）應以當年度稅課收入至少百分之五，編列債務之還本。」

    判定計算基礎（稅基）：法律明文規定是「稅課收入」，不包含「國營事業盈餘繳庫收入」或其他非稅課收入、規費等。本題的稅課收入 = 1.5 兆元 = 15,000 億元。

    計算應編列之最低債務還本額：
    \[
    15,000 \text{億元} \times 5\% = 750 \text{ 億元}
    \]

    \item 某縣政府當年度總預算及特別預算之歲出總額，各編列 110 億元及 25 億元。該縣為調節庫款收支，所舉借之未滿一年公共債務未償餘額，不得超過：
    \begin{enumerate}[label=\Alph*.]
        \item 16.5 億元
        \item 20.25 億元
        \item 33 億元
        \item 40.5 億元
    \end{enumerate}
    解答：

    根據《公共債務法》第 5 條第 10 項規定：「直轄市、縣（市）及鄉（鎮、市）為調節庫款收支所舉借之未滿一年公共債務未償餘額，不得超過其當年度總預算及特別預算歲出總額百分之三十。」

    計算歲出總額：當年度總預算（110 億元）＋ 當年度特別預算（25 億元）＝ 135 億元。
    計算未滿一年債務上限： 歲出總額 $\times$ 30\%
    \[
    135 \text{億元} \times 30\% = 40.5 \text{ 億元}
    \]


    \item 若李嘉圖等值定理（Ricardian equivalence theorem）成立，則採行減稅而改發行公債之後，國民儲蓄會有何變化？
    \begin{enumerate}[label=\Alph*.]
        \item 上揚
        \item 不變
        \item 下降
        \item 不確定
    \end{enumerate}
    解答：
    國民儲蓄 = 私人儲蓄 + 政府儲蓄（公共儲蓄）
	\begin{enumerate}[1.]
	    \item 政府儲蓄（下降）：政府減稅但支出不變，導致政府赤字增加，因此政府儲蓄減少。
        \item 私人儲蓄（上升）：依據李嘉圖等值定理，民眾預期未來會加稅來還債，所以會將現在減稅得到的錢全數存起來不消費，因此私人儲蓄增加，且增加的幅度剛好等於政府減稅的幅度。
	\end{enumerate}
    私人儲蓄增加的金額，剛好完全抵銷了政府儲蓄減少的金額。一增一減相互抵銷之下，總體經濟的國民儲蓄維持不變。

\end{enumerate}






\end{document}
